The following list shows the internal signals and per-opcode decode
behaviour of the Instruction Decoder, matching
\texttt{RWU\_RV/RV\_NoPipeline/src/asInstrDecode.sv} exactly (a single
combinational \texttt{case (instr\_opcode\_i)} statement; all outputs
default to \texttt{ALU\_ADD}/\texttt{BR\_NONE}/0/\texttt{SRC\_REGA}/
\texttt{IMM\_NONE}/\texttt{RES\_ALU} before the case, so every branch below
only lists the outputs it changes from that default).
\begin{itemize}
  \item Signal: \texttt{instr\_opcode\_i}
  \begin{itemize}
    \item Width: 7\,bit unsigned
    \item Source: Instruction from I-Mem/I-Cache.
    \item Destination: \texttt{case} statement generating all other outputs.
    \item Explanation:
    \begin{itemize}
      \item opcode=15 (\texttt{0x0F}, \texttt{fence}/\texttt{fence.i}):
      \begin{itemize}
        \item All outputs stay at their default value (no register write,
          no D-Mem access, \texttt{ALU\_ADD}/\texttt{BR\_NONE},
          \texttt{trap\_illegal\_instr\_o}=0). Decoded purely on opcode;
          \texttt{instr\_func3\_i} is not inspected.
      \end{itemize}
      \item opcode=3 (loads):
      \begin{itemize}
        \item en\_regWr\_o=1; mux\_aluSrcB\_o=1 (immediate); sel\_immSrc\_o=IMM\_I;
          mux\_resultSrc\_o=RES\_MEM; alu\_op\_o=ALU\_ADD (address = rs1+imm);
          en\_dMemRd\_o=1
      \end{itemize}
      \item opcode=19 (immediate arithmetic/logic/shift):
      \begin{itemize}
        \item en\_regWr\_o=1; mux\_aluSrcB\_o=1; sel\_immSrc\_o=IMM\_I;
          alu\_op\_o selected by \texttt{instr\_func3\_i}/\texttt{instr\_func7b5\_i}
          (see below); \texttt{trap\_illegal\_instr\_o} cannot be set here
          -- all 8 \texttt{func3} values are decoded.
      \end{itemize}
      \item opcode=23 (auipc):
      \begin{itemize}
        \item en\_regWr\_o=1; sel\_immSrc\_o=IMM\_U; mux\_aluSrcA\_o=SRC\_PC;
          mux\_aluSrcB\_o=1; alu\_op\_o=ALU\_ADD
      \end{itemize}
      \item opcode=27 (immediate arithmetic/shift, word):
      \begin{itemize}
        \item en\_regWr\_o=1; mux\_aluSrcB\_o=1; sel\_immSrc\_o=IMM\_I;
          alu\_op\_o $\in$ \{ALU\_ADDW, ALU\_SLLW, ALU\_SRLW, ALU\_SRAW\}
          (see \texttt{instr\_func3\_i} below); only \texttt{func3}
          $\in$ \{000,001,101\} are valid -- any other value sets
          \texttt{trap\_illegal\_instr\_o}=1.
      \end{itemize}
      \item opcode=35 (stores):
      \begin{itemize}
        \item en\_dMemWr\_o=1; mux\_aluSrcB\_o=1; sel\_immSrc\_o=IMM\_S;
          alu\_op\_o=ALU\_ADD
      \end{itemize}
      \item opcode=51 (R-type, RV32I arithmetic/logic/shift):
      \begin{itemize}
        \item en\_regWr\_o=1; mux\_aluSrcA\_o=SRC\_REGA; mux\_aluSrcB\_o=0
          (register); alu\_op\_o selected by \texttt{instr\_func3\_i}/
          \texttt{instr\_func7b5\_i} (see below); all 8 \texttt{func3}
          values decoded, \texttt{trap\_illegal\_instr\_o} cannot be set here.
      \end{itemize}
      \item opcode=55 (lui):
      \begin{itemize}
        \item en\_regWr\_o=1; sel\_immSrc\_o=IMM\_U; mux\_aluSrcA\_o=SRC\_ZERO;
          mux\_aluSrcB\_o=1; alu\_op\_o=ALU\_ADD (result = 0 + imm)
      \end{itemize}
      \item opcode=59 (R-type, RV64I word arithmetic/shift):
      \begin{itemize}
        \item en\_regWr\_o=1; alu\_op\_o $\in$ \{ALU\_ADDW/SUBW, ALU\_SLLW,
          ALU\_SRLW/SRAW\}; only \texttt{func3} $\in$ \{000,001,101\} are
          valid -- any other value sets \texttt{trap\_illegal\_instr\_o}=1.
      \end{itemize}
      \item opcode=99 (conditional branches):
      \begin{itemize}
        \item sel\_immSrc\_o=IMM\_B; br\_op\_o selected by
          \texttt{instr\_func3\_i} (see below); only \texttt{func3}
          $\in$ \{000,001,100,101,110,111\} are valid -- 010 and 011 are
          reserved and set \texttt{trap\_illegal\_instr\_o}=1.
      \end{itemize}
      \item opcode=103 (jalr):
      \begin{itemize}
        \item en\_regWr\_o=1; sel\_immSrc\_o=IMM\_I; mux\_jump\_o=1;
          mux\_resultSrc\_o=RES\_PC4; br\_op\_o=BR\_ALWAYS
      \end{itemize}
      \item opcode=111 (jal):
      \begin{itemize}
        \item en\_regWr\_o=1; sel\_immSrc\_o=IMM\_J; mux\_resultSrc\_o=RES\_PC4;
          br\_op\_o=BR\_ALWAYS
      \end{itemize}
      \item opcode=115 (\texttt{0x73}, SYSTEM -- CSR/\texttt{mret}/
        \texttt{ecall}/\texttt{ebreak}):
      \begin{itemize}
        \item mux\_resultSrc\_o=RES\_CSR always; en\_regWr\_o=1 unless
          \texttt{instr\_func3\_i}=000 (\texttt{mret}/\texttt{ecall}/
          \texttt{ebreak} do not write \texttt{rd}). The decoder does not
          further distinguish \texttt{mret}, \texttt{ecall}, and
          \texttt{ebreak} from one another (all three are legal,
          \texttt{func3}=000, differing only in \texttt{instr[31:20]}) --
          that distinction, and the CSR address decode for the six
          \texttt{csrr*} variants, happens outside this module in the CPU
          top level (\texttt{as\_cpux}/\texttt{cc\_cpupipefwd}); see
          Section~\ref{subsec:irqconcept}.
      \end{itemize}
      \item default (any other opcode):
      \begin{itemize}
        \item trap\_illegal\_instr\_o=1; all other outputs stay at their
          default (inert) value.
      \end{itemize}
    \end{itemize}
  \end{itemize}
  \item Signal: \texttt{instr\_func3\_i}
  \begin{itemize}
    \item Width: 3\,bit, unsigned
    \item Source: Instruction from I-Mem/I-Cache.
    \item Destination: nested \texttt{case} statements selecting
      \texttt{alu\_op\_o} (opcodes 19/27/51/59) or \texttt{br\_op\_o}
      (opcode 99); also gates \texttt{en\_regWr\_o} for opcode 115.
    \item Explanation (\texttt{alu\_op\_o} selection; \texttt{instr\_func7b5\_i}
      distinguishes the two rows sharing the same \texttt{func3}):
    \begin{itemize}
      \item opcode=19 (immediate, 64-bit width): func3=000 $\to$
        \texttt{ALU\_ADD}; 001 $\to$ \texttt{ALU\_SLL}; 010 $\to$
        \texttt{ALU\_SLT}; 011 $\to$ \texttt{ALU\_SLTU}; 100 $\to$
        \texttt{ALU\_XOR}; 101 $\to$ \texttt{ALU\_SRL}
        (\texttt{func7b5\_i}=0) or \texttt{ALU\_SRA} (=1); 110 $\to$
        \texttt{ALU\_OR}; 111 $\to$ \texttt{ALU\_AND}
      \item opcode=27 (immediate, word width): func3=000 $\to$
        \texttt{ALU\_ADDW}; 001 $\to$ \texttt{ALU\_SLLW}; 101 $\to$
        \texttt{ALU\_SRLW} (\texttt{func7b5\_i}=0) or \texttt{ALU\_SRAW} (=1)
      \item opcode=51 (register, 64-bit width): func3=000 $\to$
        \texttt{ALU\_ADD} (\texttt{func7b5\_i}=0) or \texttt{ALU\_SUB} (=1);
        001 $\to$ \texttt{ALU\_SLL}; 010 $\to$ \texttt{ALU\_SLT}; 011 $\to$
        \texttt{ALU\_SLTU}; 100 $\to$ \texttt{ALU\_XOR}; 101 $\to$
        \texttt{ALU\_SRL} (\texttt{func7b5\_i}=0) or \texttt{ALU\_SRA} (=1);
        110 $\to$ \texttt{ALU\_OR}; 111 $\to$ \texttt{ALU\_AND}
      \item opcode=59 (register, word width): func3=000 $\to$
        \texttt{ALU\_ADDW} (\texttt{func7b5\_i}=0) or \texttt{ALU\_SUBW} (=1);
        001 $\to$ \texttt{ALU\_SLLW}; 101 $\to$ \texttt{ALU\_SRLW}
        (\texttt{func7b5\_i}=0) or \texttt{ALU\_SRAW} (=1)
    \end{itemize}
    \item Explanation (\texttt{br\_op\_o} selection, opcode=99): func3=000
      $\to$ \texttt{BR\_EQ} (\texttt{beq}); 001 $\to$ \texttt{BR\_NE}
      (\texttt{bne}); 100 $\to$ \texttt{BR\_LT} (\texttt{blt}); 101 $\to$
      \texttt{BR\_GE} (\texttt{bge}); 110 $\to$ \texttt{BR\_LTU}
      (\texttt{bltu}); 111 $\to$ \texttt{BR\_GEU} (\texttt{bgeu}); 010/011
      $\to$ \texttt{trap\_illegal\_instr\_o}=1 (reserved encodings).
  \end{itemize}
  \item Signal: \texttt{instr\_func7b5\_i}
  \begin{itemize}
    \item Width: 1\,bit, unsigned
    \item Source: Instruction from I-Mem/I-Cache (instruction bit [30]).
    \item Destination: the \texttt{alu\_op\_o} case branches listed above.
    \item Explanation: distinguishes ADD/SUB and the logical/arithmetic
      shift variants (SRL/SRA, SRLI/SRAI, SRLW/SRAW, SRLIW/SRAIW) that
      share the same \texttt{func3} encoding.
  \end{itemize}
  \item Signal: \texttt{mux\_resultSrc\_o}
  \begin{itemize}
    \item Width: 2\,bit, \texttt{result\_src\_t} enum
    \item Source: generated from the opcode, see \texttt{instr\_opcode\_i} above.
    \item Destination: WB-stage result mux (ALU result / D-Mem load data /
      PC+4 / CSR read value $\to$ register file write data).
    \item Explanation: \texttt{RES\_ALU}=0 (default, most instructions),
      \texttt{RES\_MEM}=1 (loads), \texttt{RES\_PC4}=2 (\texttt{jal}/\texttt{jalr}),
      \texttt{RES\_CSR}=3 (opcode 115, SYSTEM).
  \end{itemize}
  \item Signal: \texttt{en\_dMemWr\_o}
  \begin{itemize}
    \item Width: 1\,bit, unsigned
    \item Source: generated from the opcode (stores only).
    \item Destination: D-Mem/D-Cache write enable.
    \item Explanation: Write-enable for the data memory.
  \end{itemize}
  \item Signal: \texttt{en\_dMemRd\_o}
  \begin{itemize}
    \item Width: 1\,bit, unsigned
    \item Source: generated from the opcode (loads only).
    \item Destination: D-Mem/D-Cache read enable; also gates the load-use
      hazard stall in \texttt{RV\_PipelineCacheFwd} (Section~\ref{subsec:pipehazard}).
    \item Explanation: Read-enable for the data memory.
  \end{itemize}
  \item Signal: \texttt{mux\_aluSrcB\_o}
  \begin{itemize}
    \item Width: 1\,bit, unsigned
    \item Source: generated from the opcode.
    \item Destination: ALU source-B input mux.
    \item Explanation: \[ \begin{array}{ll}
                                              0 \Rightarrow & \mbox{register (rs2) to ALU source B} \\
                                              1 \Rightarrow & \mbox{immediate value to ALU source B}
                                              \end{array} \]
  \end{itemize}
  \item Signal: \texttt{mux\_aluSrcA\_o}
  \begin{itemize}
    \item Width: 2\,bit, \texttt{mux\_a\_t} enum
    \item Source: generated from the opcode.
    \item Destination: ALU source-A input mux.
    \item Explanation: \texttt{SRC\_REGA}=0 (register rs1, default),
      \texttt{SRC\_PC}=1 (\texttt{auipc}), \texttt{SRC\_ZERO}=2 (\texttt{lui}).
  \end{itemize}
  \item Signal: \texttt{en\_regWr\_o}
  \begin{itemize}
    \item Width: 1\,bit, unsigned
    \item Source: generated from the opcode (and, for opcode 115, \texttt{instr\_func3\_i}).
    \item Destination: Register-file write enable.
    \item Explanation: Write enable for the register file.
  \end{itemize}
  \item Signal: \texttt{mux\_jump\_o}
  \begin{itemize}
    \item Width: 1\,bit, unsigned
    \item Source: opcode=103 (\texttt{jalr}) only.
    \item Destination: Branch-target-adder source-A select.
    \item Explanation: \[ \begin{array}{ll}
                                              0 \Rightarrow & \mbox{PC to branch-target-adder input A (default: branches, JAL)} \\
                                              1 \Rightarrow & \mbox{register rs1 to branch-target-adder input A (JALR)}
                                              \end{array} \]
  \end{itemize}
  \item Signal: \texttt{sel\_immSrc\_o}
  \begin{itemize}
    \item Width: 3\,bit, \texttt{imm\_src\_t} enum
    \item Source: generated from the opcode.
    \item Destination: Immediate generator format select.
    \item Explanation: \[ \begin{array}{ll}
                                              \texttt{IMM\_NONE}=0 \Rightarrow & \mbox{no immediate (R-type, FENCE, default)} \\
                                              \texttt{IMM\_I}=1 \Rightarrow & \mbox{I-type (loads, immediate arithmetic, jalr)} \\
                                              \texttt{IMM\_S}=2 \Rightarrow & \mbox{S-type (stores)} \\
                                              \texttt{IMM\_B}=3 \Rightarrow & \mbox{B-type (branches)} \\
                                              \texttt{IMM\_U}=4 \Rightarrow & \mbox{U-type (lui, auipc)} \\
                                              \texttt{IMM\_J}=5 \Rightarrow & \mbox{J-type (jal)}
                                              \end{array} \]
  \end{itemize}
  \item Signal: \texttt{alu\_op\_o}
  \begin{itemize}
    \item Width: 6\,bit declared (\texttt{alu\_op\_t}), 15 values used
    \item Source: generated from the opcode and, where applicable,
      \texttt{instr\_func3\_i}/\texttt{instr\_func7b5\_i} (see above).
    \item Destination: Algorithmic ALU operation select.
    \item Explanation: \[ \begin{array}{ll}
                                              0 \Rightarrow & \texttt{ALU\_ADD} \mbox{ (ADD, ADDI, loads, stores, jal/jalr target, auipc, lui, default)} \\
                                              1 \Rightarrow & \texttt{ALU\_SUB} \\
                                              2 \Rightarrow & \texttt{ALU\_AND} \\
                                              3 \Rightarrow & \texttt{ALU\_OR} \\
                                              4 \Rightarrow & \texttt{ALU\_XOR} \\
                                              5 \Rightarrow & \texttt{ALU\_SLT} \\
                                              6 \Rightarrow & \texttt{ALU\_SLTU} \\
                                              7 \Rightarrow & \texttt{ALU\_SLL} \\
                                              8 \Rightarrow & \texttt{ALU\_SRL} \\
                                              9 \Rightarrow & \texttt{ALU\_SRA} \\
                                              10 \Rightarrow & \texttt{ALU\_ADDW} \\
                                              11 \Rightarrow & \texttt{ALU\_SUBW} \\
                                              12 \Rightarrow & \texttt{ALU\_SLLW} \\
                                              13 \Rightarrow & \texttt{ALU\_SRLW} \\
                                              14 \Rightarrow & \texttt{ALU\_SRAW}
                                              \end{array} \]
  \end{itemize}
  \item Signal: \texttt{br\_op\_o}
  \begin{itemize}
    \item Width: 3\,bit, \texttt{br\_op\_t} enum
    \item Source: opcode 99 (branches, via \texttt{instr\_func3\_i}) and
      opcodes 103/111 (\texttt{jalr}/\texttt{jal}, unconditional).
    \item Destination: Branch ALU operation select (\texttt{take\_o}
      generation, both chip variants -- see Section~\ref{sec:alufunc01}).
    \item Explanation: \[ \begin{array}{ll}
                                              0 \Rightarrow & \texttt{BR\_NONE} \mbox{ (default, no branch/jump)} \\
                                              1 \Rightarrow & \texttt{BR\_EQ} \mbox{ (beq)} \\
                                              2 \Rightarrow & \texttt{BR\_NE} \mbox{ (bne)} \\
                                              3 \Rightarrow & \texttt{BR\_LT} \mbox{ (blt)} \\
                                              4 \Rightarrow & \texttt{BR\_GE} \mbox{ (bge)} \\
                                              5 \Rightarrow & \texttt{BR\_LTU} \mbox{ (bltu)} \\
                                              6 \Rightarrow & \texttt{BR\_GEU} \mbox{ (bgeu)} \\
                                              7 \Rightarrow & \texttt{BR\_ALWAYS} \mbox{ (jal, jalr -- unconditional)}
                                              \end{array} \]
  \end{itemize}
  \item Signal: \texttt{trap\_illegal\_instr\_o}
  \begin{itemize}
    \item Width: 1\,bit, unsigned
    \item Source: the case-statement default (unrecognized opcode) and the
      \texttt{default} arms of the \texttt{instr\_func3\_i} sub-cases for
      opcodes 27, 59, and 99 (the only opcodes whose 8 possible
      \texttt{func3} values are not all valid RV64I encodings).
    \item Destination: CPU top-level trap logic (\texttt{as\_cpux}/
      \texttt{cc\_cpupipefwd}); sets \texttt{mcause}=2.
    \item Explanation: see Section~\ref{subsec:irqconcept} for the full
      trap-priority chain and CSR update sequence.
  \end{itemize}
\end{itemize}
